\documentclass[10pt,a4paper]{amsart}
\usepackage[margin=19mm]{geometry}
\usepackage{amsmath,amssymb,mathtools}
\usepackage[scr=rsfs,cal=euler]{mathalfa}
\usepackage{xfrac}
\usepackage[unicode,hidelinks,bookmarks=false]{hyperref}
\newcommand{\ie}{i.e.\ }
\newcommand{\cf}{cf.\ }
\newcommand{\resp}{respectively\ }
\newcommand{\Subsection}{\subsection}
\providecommand{\nobreakdash}{\nobreak\hbox{-}}
\setlength{\emergencystretch}{2em}
\multlinegap0pt
\usepackage{bm}
\usepackage{bbm}
\usepackage{dsfont}
\usepackage{xspace}
\usepackage{cancel}
\usepackage{mathtools}
\usepackage{amsthm}
\makeatletter
\@ifundefined{lemma}{\newtheorem{lemma}{Lemma}}{}
\makeatother
\def\a{\alpha}
\def\b{\beta}
\def\g{\gamma}
\title[Spectral properties of Schr\"odinger operator with translati]{Spectral properties of Schr\"odinger operator with translations and Neumann boundary conditions}
\author{D.I. Borisov, D.M. Polyakov}
\date{Source: arXiv:2506.23645v1 (CC BY 4.0)}
\begin{document}
\maketitle

\noindent\emph{Source and licence.} Spectral properties of Schr\"odinger operator with translations and Neumann boundary conditions. Version arXiv:2506.23645v1 (CC BY 4.0), distributed under the Creative Commons Attribution 4.0 International licence, as recorded in the arXiv licence metadata for this exact version and shown on the abstract page https://arxiv.org/abs/2506.23645v1. Full rights notice: licenses/arxiv-2506.23645-CC-BY-4.0.txt.\par\smallskip

\noindent\textit{Editorial scope.} Counted unit: the Lemma whose source label is \texttt{lm4.3} in the article (source0.tex lines 883--893, source label \texttt{lm4.3}), delivered together with the article's own complete original proof of it (source0.tex lines 894--980, one \texttt{proof} environment of 87 lines). No mathematics is written, repaired or re-derived: statement and proof are the article's own text line for line. Required context retained uncounted: equateta (lines 845--847); eqmaineta (lines 849--852); 4.16 (lines 859--862); estAgamma (lines 864--866); varphiseries (lines 877--879). Every definition, display and label of those lines is retained unchanged. Omitted from this package and not redistributed: 1--844, 848--848, 853--858, 863--863, 867--876, 880--882, 981--1487. Nothing inside a retained excerpt is omitted or altered: every retained body is delivered exactly as the article's lines are written, so no omission receipt of any kind accompanies this package. Citations: the retained excerpts carry no citation command at all, so no bibliography entry of the article is transcribed and no reference is redistributed. Reference adaptation: none was needed and none was made; every reference inside the delivered unit resolves inside it, so no unresolved reference or citation is produced. Printed slips kept literal: no printed word, symbol, hypothesis or number of the counted statement or of its proof was altered. Layout and class: the article's own class is \texttt{article} with packages including fontenc, babel, amssymb, amsmath, amsfonts, amsthm, color, graphicx, srcltx, dsfont, geometry; the wrapper is the lane's pinned \texttt{amsart} editorial wrapper at 10pt on A4, which restates the theorem-like environments the retained text needs and adds the article's own macro definitions that the retained text needs (\texttt{\textbackslash a}, \texttt{\textbackslash b}, \texttt{\textbackslash g}), with extra wrapper packages (bm, bbm, dsfont, xspace, cancel, mathtools). The delivered numbering is the unit's own. Assets: no retained span contains a figure environment, a graphics-inclusion command, a PDF-page inclusion, a table or a TikZ picture; no image file is copied into this package, the package declares no asset dependency and no image file is redistributed with it. Licence: version arXiv:2506.23645v1 of this article is distributed under the Creative Commons Attribution 4.0 International licence, as recorded in the arXiv licence metadata for this exact version and shown on the abstract page https://arxiv.org/abs/2506.23645v1. A full rights notice is delivered at licenses/arxiv-2506.23645-CC-BY-4.0.txt. Characters: every retained excerpt is pure ASCII, so the delivered engine, XeLaTeX with its default Latin Modern text font, reports no missing character.
\begin{equation}\label{equateta}
\eta-\mu F(\eta, \mu, n,\a,\b)=0,
\end{equation}

\begin{equation}\label{eqmaineta}
F(\eta, \mu, n,\a,\b)=\sum\limits_{j=0}^{\infty}
\frac{\mu^j f_j (\pi n+\arcsin\eta,\a,\b)}{(1+\mu\arcsin\eta)^{j+1}}.
\end{equation}

\begin{equation}\label{4.16}
F (\eta, \mu, n,\a,b)=\sum\limits_{\gamma\in\mathds{Z}_+^2}A_\gamma(n,\a,\b)\eta^{\gamma_1}
\mu^{\gamma_2}, \qquad \gamma=(\gamma_1, \gamma_2),
\end{equation}

\begin{equation}\label{estAgamma}
|A_\gamma(n,\a,\b)|\leqslant c_{8}c_{9}^{|\gamma|},
\end{equation}

\begin{equation}\label{varphiseries}
\varphi(\mu, n,\a,\b)=\mu\sum\limits_{i=0}^\infty B_i(n,\a,\b)\mu^i,
\end{equation}

\begin{lemma}\label{lm4.3}
There exists a fixed positive constant $c_{10}$   independent of $\alpha$, $\beta$, and $n$ such that
$$
r(n,\a,\b)\geqslant c_{10}>0\quad \text{for all}\quad n.
$$
The coefficients $B_i$, $i\geqslant 0$, satisfy the estimate
\begin{equation}\label{4.24**}
|B_i(n,\a,\b)|\leqslant c_{11} c_{12}^i,\qquad i\geqslant 0,
\end{equation}
where $c_{11}$, $c_{12}$ are positive constants independent of $\alpha$, $\beta$, and $n$.
\end{lemma}

\begin{proof}
According to the Cauchy--Hadamard formula, it is sufficient to
estimate the coefficients $B_i$ uniformly in $n$. In order to do it, we use appropriate majorants for the functions
$F$ and $\varphi$. As the majorant for the former, we employ the holomorphic function
$\widetilde{F}=\widetilde{F}(\eta,\mu)$ defined
by the right hand side of~(\ref{4.16}) with the coefficients $A_\g$ replaced by $c_{8} c_{9}^{|\g|},$ namely,
\begin{equation*}
\widetilde{F}(\eta,\mu):=\sum\limits_{\gamma\in\mathds{Z}_+^2}
c_{8} c_{9}^{|\g|} \eta^{\gamma_1}\mu^{\gamma_2}=\frac{c_{8}}{(1-c_{9}\eta)(1-c_{9}\mu)}.
\end{equation*}
Then we consider   Equation \eqref{equateta} with such a function, and we rewrite this equation as
$$
c_{9}\eta^2(\mu)-\eta(\mu)+\frac{c_{8}\mu}{1-c_{9}\mu}=0.
$$
This equation has the unique solution, which vanishes at $\mu=0,$  and this solution can be found explicitly
$$
\eta(\mu)=\frac{1-\sqrt{1-\frac{4c_{8}c_{9}\mu}{1-c_{9}\mu} }}{2c_{9}}.
$$
It is straightforward to verify that the coefficients of    Taylor series about $\mu=0$ of the functions
\begin{equation*}
\mu\mapsto \frac{4c_{8}c_{9}\mu}{1-c_{9}\mu},\qquad u\mapsto \frac{1-\sqrt{1-u}}{2c_{9}}
\end{equation*}
are real and positive. Hence, the same is true for the function $\eta(\mu)$ as well as for
the function
$$
\widetilde{\varphi}(\mu):=\frac{1-\sqrt{1-\frac{4c_{8}c_{9}\mu}{1-c_{9}\mu}}}{2c_{9}\mu}.
$$
It is clear that the latter function is holomorphic in sufficiently small $\mu$ and by $\tau_i,$ $i\geqslant 0$, we denote its
Taylor coefficients,
\begin{equation*}
\widetilde{\varphi}(\mu)=\sum\limits_{i=0}^{\infty} \tau_i\mu^i,\qquad \tau_0=c_{8},\qquad \tau_i>0.
\end{equation*}
Since the function $\widetilde{\varphi}$ is holomorphic, we have standard bounds for $\tau_i$
\begin{equation*}%\label{4.23}
0<\tau_i\leqslant c_{11} c_{12}^i,
\end{equation*}
where $c_{11}$ and $c_{12}$ are some positive constants independent of $i,$ and of course, of $\alpha$, $\beta$, and $n$.

Now we are going to show that
\begin{equation}\label{4.24}
|B_i(n,\a,\b)|\leqslant \tau_i,\qquad i\geqslant0.
\end{equation}
We substitute the representation \eqref{varphiseries} into \eqref{eqmaineta} and obtain
%\begin{equation}\label{4.18}
\begin{align*}
\sum\limits_{i=0}^\infty B_i(n,\a,\b)\mu^i=&\sum\limits_{\gamma\in\mathds{Z}_+^2}
A_\gamma(n,\a,\b)\mu^{\gamma_2}
\bigg(\sum\limits_{i=0}^\infty B_i(n,\a,\b)\mu^i\bigg)^{\gamma_1}
\\
=&\sum\limits_{i=0}^\infty\mu^i\sum\limits_{\gamma\in\mathds{Z}_+^2}
\sum\limits_{\substack{\varkappa\in\mathds{Z}_+^{\gamma_1} \\ |\varkappa|+\gamma_2=i}}
A_\gamma(n,\a,\b)B_{\varkappa_1}(n,\a,\b)\cdot\ldots\cdot B_{\varkappa_{\gamma_1}}(n,\a,\b),
\end{align*}
%\end{equation}
where $\varkappa=(\varkappa_1, \dots, \varkappa_{\gamma_1})$. Hence, the coefficients
$B_i$ satisfy the recurrent relations
\begin{equation}\label{maincoeff}
B_i(n,\a,\b)=\sum\limits_{\gamma\in\mathds{Z}_+^2}
\sum\limits_{\substack{\varkappa\in\mathds{Z}_+^{\gamma_1} \\ |\varkappa|+\gamma_2=i}}
A_\gamma(n,\a,\b)B_{\varkappa_1}(n,\a,\b)\cdot\ldots\cdot B_{\varkappa_{\gamma_1}}(n,\a,\b),
\end{equation}
where we adopt that $B_{\varkappa_1}(n,\a,\b)\cdot\dots\cdot B_{\varkappa_{\gamma_1}}(n,\a,\b):=1$
if $\gamma_1=\gamma_2=0$ or $\varkappa=0$.
Then the formula \eqref{maincoeff} expresses the coefficients $B_i$ in terms of $B_0, \dots, B_{i-1}$.

It is clear that $B_0(n,\a,\b)=A_0(n,\a,\b)$ and hence, by \eqref{estAgamma}, we immediately get
the estimate \eqref{4.24} for $B_0(n,\a,\b)$. We proceed by induction. Suppose that the estimate
\eqref{4.24} is valid for  $i\leqslant k-1$ and let us establish it for $i=k$. Using the induction assumption,
the formulas \eqref{estAgamma}, \eqref{maincoeff}, and the positivity of the coefficients $\tau_i$, we obtain
\begin{equation}\label{prom1''}
\begin{aligned}
|B_k(n,\a,\b)|\leqslant & \sum\limits_{\gamma\in\mathds{Z}_+^2}
\sum\limits_{\substack{\varkappa\in\mathds{Z}_+^{\gamma_1} \\ |\varkappa|+\gamma_2=k}}
|A_\gamma(n,\a,\b)||B_{\varkappa_1}(n,\a,\b)|\cdot\dots\cdot |B_{\varkappa_{\gamma_1}}(n,\a,\b)|
\\
\leqslant & c_{8}\sum\limits_{\gamma\in\mathds{Z}_+^2}
\sum\limits_{\substack{\varkappa\in\mathds{Z}_+^{\gamma_1} \\ |\varkappa|+\gamma_2=k}}
c_{9}^{|\gamma|}\tau_{\varkappa_1}\cdot\dots\cdot\tau_{\varkappa_{\gamma_1}}.
\end{aligned}
\end{equation}
Then we make the same calculations with the function
$\widetilde{F}(\mu\widetilde{\varphi}(\mu), \mu)$ and we see easily that the coefficients at $\mu^k$ in its Taylor series is exactly the right hand side of \eqref{prom1''}. Since the function $\mu\widetilde{\varphi}(\mu)$ solves Equation~\eqref{equateta}
with $F$ replaced by $\widetilde{F}$, we conclude that the coefficient at $\mu^k$ in the Taylor series for $\widetilde{F}(\mu\widetilde{\varphi}(\mu), \mu)$ coincides with the coefficients at $\mu^k$ in the Taylor series for $\widetilde{\varphi}(\mu)$, that is, with $\tau_k$. Then we can replace the right hand side of \eqref{prom1''} by $\tau_k$ and this proves the estimate \eqref{4.24**}
for $B_k(n, \alpha, \beta)$.

Now we apply the Cauchy--Hadamard formula for the convergence radius of the series \eqref{varphiseries} and see that it is estimated from below by the constant $c_{10}$. The proof is complete.
\end{proof}

\end{document}
